\section{Conclusion}


FHE and MPC solve complementary halves of the private computation problem. FHE eliminates the need to share data; MPC eliminates the single point of trust for the decryption key. The hybrid construction---FHE computation with FROST-based threshold decryption---provides a practical framework for multi-party private computation on blockchain networks.

The implementation on the Lux network (F-Chain for FHE, M-Chain for MPC) achieves 94ms end-to-end latency for private orderbook matching, with configurable trust thresholds and no party ever observing the full order book in cleartext.

\begin{thebibliography}{9}
\bibitem{cggi2020} Chillotti, Gama, Georgieva, Izabach\`{e}ne, ``TFHE: Fast Fully Homomorphic Encryption over the Torus,'' Journal of Cryptology, 2020.
\bibitem{spdz2012} Damg\r{a}rd, Pastro, Smart, Zakarias, ``Multiparty Computation from Somewhat Homomorphic Encryption,'' CRYPTO 2012.
\bibitem{frost2020} Komlo, Goldberg, ``FROST: Flexible Round-Optimized Schnorr Threshold Signatures,'' SAC 2020.
\bibitem{lux-tfhe} Lux Industries, Inc., ``Torus Threshold Fully Homomorphic Encryption: Boolean Circuits on Encrypted Data with Distributed Decryption,'' 2025.
\bibitem{lux-mpc} Lux Industries, Inc., ``M-Chain: Multi-Party Computation on the Lux Network,'' 2025.
\bibitem{shamir1979} Shamir, ``How to Share a Secret,'' Communications of the ACM, 1979.
\end{thebibliography}

